\documentclass[5pt, a4paper]{article}

\usepackage[utf8]{inputenc}
\usepackage[T1]{fontenc}
\usepackage{textcomp}
\usepackage{amsmath, amssymb}
\usepackage{multicol}
\usepackage{proof}
\usepackage{enumitem}
\usepackage{fancyhdr}
\usepackage[margin=1in]{geometry}
% figure support
\usepackage{import}
\usepackage{xifthen}
\pdfminorversion=7

\usepackage[most]{tcolorbox}
\usepackage{pdfpages}
\usepackage{transparent}
\newcommand{\incfig}[1]{%
    \def\svgwidth{\columnwidth}
    \import{./figures/}{#1.pdf_tex}
}

\tcbset{
colback=gray!5,
colframe=blue!50!white,
coltitle=black,
arc=0pt,
boxrule=-0.1mm,
fonttitle=\bfseries,
coltitle=black,
left=2.5mm,
leftrule=1mm,
right=3.5mm,
colbacktitle=black!10!white,
toptitle=0.75mm,
bottomtitle=0.5mm,
}
\begin{document}
\begin{tcolorbox}[title=Question 1]
    Find an interpolating polynomial for the data points $(0,1),(2,2), (3,4)$
\end{tcolorbox}

Solving a second power polynomial: 
\begin{align*}
    0 + 0 + c &= 1 \\
    4a + 2b + c &= 2 \\
    9a + 3b + c &= 4 
.\end{align*}
\[
    c = 1, b = -\frac{1}{2}, c = \frac{1}{2}
.\] 
\[
p(x) = \frac{1}{2}x^2 - \frac{1}{2}x + 1
,\] 
or one could interpolate using the monomial basis, 
\begin{align*}
    \begin{pmatrix}
        1 & 0 & 0 \\ 
        1 & 2 & 4 \\ 
        1 & 3 & 9 \\ 
    \end{pmatrix} 
    \begin{pmatrix} a \\ b \\ c \end{pmatrix} 
    &= \begin{pmatrix} 1 \\ 2 \\ 4 \end{pmatrix}  
.\end{align*}
Which would yield the same result.


\begin{tcolorbox}[title=Question 2]
    Consider the polynomial interpolation of the function $f(x)$ with nodes  \[
    x_0=-3, x_1=-2, x_2 = -1, x_3 = 0, x_4=1, x_5 = 2, x_6 = 3
    .\] 
    Suppose that the interpolating polynomial is quadratic. Given that \[
    f(x_0) = -9, f(x_4) = 3, f(x_5) = -4
    ,\] 
    find the values of $f(x)$ on the other interpolation nodes.
\end{tcolorbox}

Since we are given three data points, the solution for $f(x)$ would be unique,
and the values of $f(x)$ on the other nodes would be as well. Using a monomial
basis,
\begin{align*}
    \begin{pmatrix} 
        1 & -3 & 9 \\ 
        1 & 1 & 1 \\ 
        1 & 2 & 4 \\
    \end{pmatrix} 
    \begin{pmatrix} a \\ b \\ c \end{pmatrix} 
    &= \begin{pmatrix} -9 \\ 3 \\ -4 \end{pmatrix}  
.\end{align*}
\[
f(x) = 6x^2 + x - 2
.\] 
Then, subtitute the values of the nodes into the function and we'll arrive at
the answer.

\begin{tcolorbox}[title=Question 3]
    Given 2 disticnt points $x_0=0, x_1=1$ and that \[
    f(0) = f_0, f(1) = f_1, f'(0) =  f_0', f'(1) = f_1'
    .\] 
    where $f'(x_{i})$ is the derivative of $f$ at $x_{i}$.
    Write a polynomial $P$ with three degrees in terms of  $f_0,f_1,f_0',f_1'$
    such that \[
    P(x_{i}) = f(x_{i}), P'(x_{i}) = f'(x_{i}), \forall i = 0,1
    .\] 
\end{tcolorbox}
Solve for the original function \[
    \begin{pmatrix} 
        1 & 0 & 0 & 0 \\
        1 & 1 & 1 & 1\\
\end{pmatrix} 
\begin{pmatrix} a \\ b \\ c \\ d \end{pmatrix} 
= \begin{pmatrix} f_0 \\ f_1 \end{pmatrix} 
.\] 
and for the derivatives \[
    \begin{pmatrix} 
        0 & 0 & 0 \\
        1 & 1 & 1\\
    \end{pmatrix} 
    \begin{pmatrix} b \\ 2c \\ 3d \end{pmatrix} 
    =  \begin{pmatrix} f_0' \\ f_1' \end{pmatrix} 
.\] 

Solving for $\begin{pmatrix} a & b & c & d \end{pmatrix}^T$ yields \[
a = f_0, \quad b = 2(f_1 - f_0) - f_1' - d, \quad  c = f_1' + f_0 - f_1 + 2d,
\quad d \in \mathbb{R}
.\] 
\[
P(x) = dx^3 + ( f_1' + f_0 - f_1 + 2d) + x^2 (2(f_1 - f_0) - f_1' - d) x + f_0
.\] 
\begin{tcolorbox}[title=Question 4]
    GIven the data points $(-1,2),(0,3),(1,6)$, determine tha interpolating
    polynomial of degree two:
     \begin{enumerate}
        \item using the monomial basis
        \item using the Langrage basis
        \item using the Newton basis
        \item Verify that the three methods give the same polynomial.
    \end{enumerate}
\end{tcolorbox}
Attempt using a monomial basis: 
\[
\begin{pmatrix} 
    1 & -1 & 1 \\ 
    1 & 0 & 0 \\ 
    1 & 1 & 1 \\
\end{pmatrix} 
\begin{pmatrix} a \\ b \\ c \end{pmatrix} 
= \begin{pmatrix} 2 \\ 3 \\ 6 \end{pmatrix} 
.\] 
\[
    a = 3, b = 2, c = 1
.\] 
\[
p(x) = x^2 + 2x + 3
.\] 
For a langrage basis, 
\begin{align*}
    L_0(x) &= \frac{x(x - 1)}{2} \\
    L_1(x) &=  \frac{(x+1)(x-1)}{-1}\\
    L_2(x) &= \frac{(x + 1)x}{2} 
.\end{align*}

\[
P(x) = x(x-1) - 3(x+1)(x-1) + 3(x+1)x
.\] 

\[
p(x) = x^2 + 2x + 3
.\] 
For a Newton basis:
\[
\begin{pmatrix} 
    1 & 0 & 0 \\ 
    1 & 1 & 0 \\ 
    1 & 2 & 2 \\
\end{pmatrix} 
\begin{pmatrix} a \\ b \\ c \end{pmatrix} 
= \begin{pmatrix} 2 \\ 3 \\ 6 \end{pmatrix} 
.\] 
\[
p(x) = x^2 + 2x + 3
.\] 
As you can see, the three methods yields the same polynomial.
\begin{tcolorbox}[title=Question 5]
    Consider the data $(-1,-10),(0,-4)(1,0),(2,14)$.
     \begin{enumerate}
        \item Set up and solve the normal equation to find the least-square line
            of best fit.
        \item Set up and solve the normal equation to find the least-square
            parabola of best fit.
        \item Find the interpolating cubic polynomial using the Langrage basis
            functions.
        \item Find the cubid Newton polynomial interpolant using divided
            differences.
    \end{enumerate}
\end{tcolorbox}
For least square line:
\begin{align*}
    \begin{pmatrix} 
        1 & -1 \\
        1 & 0 \\
        1 & 1 \\
        1 & 2 \\
    \end{pmatrix} 
    \begin{pmatrix} a_0 \\ a_1 \end{pmatrix} 
    &= \begin{pmatrix} -10 \\ -4 \\ 0 \\ 14 \end{pmatrix}  \\
    \begin{pmatrix} 
        1 & -1 \\
        1 & 0 \\
        1 & 1 \\
        1 & 2 \\
    \end{pmatrix}^T
    \begin{pmatrix} 
        1 & -1 \\
        1 & 0 \\
        1 & 1 \\
        1 & 2 \\
    \end{pmatrix} 
    \begin{pmatrix} a_0 \\ a_1 \end{pmatrix} 
    &= 
    \begin{pmatrix} 
        1 & -1 \\
        1 & 0 \\
        1 & 1 \\
        1 & 2 \\
    \end{pmatrix}^T
    \begin{pmatrix} -10 \\ -4 \\ 0 \\ 14 \end{pmatrix} \\
.\end{align*}
For least square parabola
\begin{align*}
    \begin{pmatrix} 
        1 & -1 & 1\\ 
        1 & 0 & 0 \\
        1 & 1 & 1 \\ 
        1 & 2 & 4 
    \end{pmatrix} 
    \begin{pmatrix} a_0 \\ a_1 \\ a_2 \end{pmatrix} 
    &= \begin{pmatrix} -10 \\ -4 \\ 0 \\ 14 \end{pmatrix}  \\
    \begin{pmatrix} 
        1 & -1 & 1 \\ 
        1 & 0 & 0 \\
        1 & 1 & 1 \\ 
        1 & 2 & 4 
    \end{pmatrix}^T
    \begin{pmatrix} 
        1 & -1 & 1 \\ 
        1 & 0 & 0 \\
        1 & 1 & 1 \\ 
        1 & 2 & 4 
    \end{pmatrix} 
    \begin{pmatrix} a_0 \\ a_1 \\ a_2 \end{pmatrix} 
    &=
    \begin{pmatrix} 
        1 & -1 & 1 \\ 
        1 & 0 & 0 \\
        1 & 1 & 1 \\ 
        1 & 2 & 4 
    \end{pmatrix}^T
    \begin{pmatrix} -10 \\ -4 \\ 0 \\ 14 \end{pmatrix}  \\
.\end{align*}
Using the Langrange basis for a cubic polynomial,
\begin{align*}
    L_0(x) &= \frac{x(x-1)(x-2)}{(-1)(-1-1)(-1-2)}
           = \frac{x(x-1)(x-2)}{-6} \\
    L_1(x) &= \frac{(x+1)(x-1)(x-2)}{(1)(-1)(-1)}
           =  (x+1)(x-1)(x-2)\\
    L_2(x) &= \frac{(x+1)(x)(x-2)}{-2}\\
    L_3(x) &= \frac{(x+1)(x)(x-1)}{6} 
.\end{align*}
Then,
\begin{align*}
    P(x) &= \frac{1}{6}x(x-1)(x-2) - \frac{1}{2}x(x+1)(x-2) +
    \frac{1}{3}x(x+1)(x-1)\\
.\end{align*}

\begin{tcolorbox}[title=Question 6]
    COnsider the data points $(-1,4),(0,1),(1,0),(2,-5)$ and the problem of
    finding an interpolant polynomial that fits these points.
     \begin{enumerate}
        \item Write down, in matrix form, the system of linear equation that
            corresponds to solving this problem with the monomial basis. You do
            not need to solve this system.
        \item Wirte down the associated Newton basis functions and determine the
            coeffecients needed to represent the interpolating polynomial in
            terms of this basis.
        \item Find a cubic polynomial that passes through the first two data
            points and has the property that $p''(-1) = 0$ and  $p'(0) = -1$
    \end{enumerate}
\end{tcolorbox}

Let $x_0 = -1, x_1 = 0, x_2 = 1, x_3 = 2$. Likewise, $y_0 = 4, y_1 = 1, y_2 = 0,
y_3 = -5$. We'll write the system in terms of the monomial basis as such:

\[
\begin{pmatrix} 
    1 & x_0 & x_0^2 \\ 
    1 & x_1 & x_1^2 \\ 
    1 & x_2 & x_2^2 \\
    1 & x_3 & x_3^2 \\
\end{pmatrix} 
 \begin{pmatrix} a_0 \\ a_1 \\ a_2 \\ a_3 \end{pmatrix} 
 = \begin{pmatrix} y_0 \\ y_1 \\ y_2 \\ y_3 \end{pmatrix} 
.\] 
Solving for $\begin{pmatrix} a_0 & a_1 & a_2 & a_3 \end{pmatrix}^T$ would
give the coeefecient in the interpolant polynomial.



\begin{tcolorbox}[title=Question 7]
    Suppose that you want to interpolate the points \[
        (-1,0),(0,1),(2,0),(3,1),(4,2)
    \] 
    by a polynomial of as low a degree as possible. What is the maximum
    neccessary degree of this polynomial? Write out a linear system of equations
    for the coeffecients of the interpolating polynomial.
\end{tcolorbox}

Since we are given 5 data points, we can draw a rough bound for the degrees
neccessary: no more than 5. We can check the actual degrees neccessary by
forming the monomial matrix and checking the rank of it's RREF.
\end{document}
